\section{Simulation results and output coupling}

\subsection{General-order verification and model consistency}

The 84-case matrix covered open and resistively loaded branches, several pairing rules, nonuniform same-transfer realizations, fixed-load matching families, reactive loads, source impedance, output coupling and two-load recovery. Orders ranged from one to twelve. With the documented 1~m$\Omega$ inductor series resistance included in the independent model, the largest absolute complex node-voltage discrepancy was $1.68\times10^{-14}$~V.

Against the originally intended ideal-inductor reference, two cases exceeded the unchanged $10^{-7}$~V gate, with errors of $4.54\times10^{-7}$ and $3.57\times10^{-7}$~V. Fresh runs specifying $R_{\rm ser}=0$ agreed with the ideal reference to less than $1.8\times10^{-15}$~V; explicit 1~m$\Omega$ runs reproduced the default-model behavior. The discrepancy arose from different element models. Both original failures and the diagnostic records were retained.

At $10$~kHz, the three-stage same-$A$/different-$B$ construction agreed in the open circuit to a relative difference of $1.4\times10^{-14}$ and separated under a 10~k$\Omega$ load by $3.4\times10^{-4}$. A fixed-load family agreed at $g=0.1$ to $6.7\times10^{-14}$ and differed at $g=0.2$ by $3.2\times10^{-5}$. A capacitive-load matching family agreed under the specified capacitance to $5.3\times10^{-14}$ and differed under resistive loading by $2.2\times10^{-3}$. The comparisons tested complete complex outputs; the following second-order examples make the same mechanisms visible with shorter component tuples.

\subsection{Termination-dependent equivalence}

Table~\ref{tab:branches} collects the three branch constructions. The equalities and inequalities were established algebraically before the frequency sweeps.

\begin{table}[htbp]
\centering
\caption{Normalized component values and branch coefficients.}
\label{tab:branches}
\begin{tabular}{lrrrrrrrr}
\toprule
Branch & $r_1$ & $r_2$ & $q_1$ & $q_2$ & $a$ & $b$ & $c$ & $d$ \\
\midrule
P & 1 & 1 & 1 & 1 & 3 & 1 & 2 & 1 \\
Q & 1 & 4 & $1/2$ & $1/2$ & 3 & 1 & 5 & 2 \\
T & $18/19$ & $20/19$ & $361/180$ & $1/2$ & 2.9 & 1 & 2 & 2 \\
\bottomrule
\end{tabular}
\end{table}

Figure~\ref{fig:termination-transfer} compares the open and loaded P/Q responses. Figure~\ref{fig:output-impedance} shows the output impedances that distinguish the open-circuit-equivalent branches. Figure~\ref{fig:load-specific} compares T with P after each load is applied.



The simulations reproduced open-circuit cancellation for P/Q and cancellation at \(R_L=10\,\mathrm{k\Omega}\) for P/T. Changing \(R_L\) to 5 k\(\Omega\) produced unequal P/T responses. Across the 37 cases, the maximum absolute complex node-voltage difference between LTspice and the independently assembled nodal solution was \(1.14\times10^{-15}\) V.

\begin{figure}[htbp]
\centering
\includegraphics[width=0.67\linewidth]{figures/Fig_03_termination_transfer.pdf}
\caption{Magnitude responses of P and Q under open-circuit and 10~k$\Omega$ shunt-load conditions. Lines are analytical transfer functions; open circles are archived LTspice samples. Open-circuit responses coincide. The normalization is $R_*=1$~k$\Omega$, $C_*=10$~nF.}
\label{fig:termination-transfer}
\end{figure}

\begin{figure}[htbp]
\centering
\includegraphics[width=0.67\linewidth]{figures/Fig_04_output_impedance.pdf}
\caption{Open-circuit output-impedance magnitudes of P, Q and T with the ideal input voltage sources set to zero. Curves are calculated from the four-coefficient representation. P and Q share an open-circuit voltage transfer but have different output impedances.}
\label{fig:output-impedance}
\end{figure}

\begin{figure}[htbp]
\centering
\includegraphics[width=0.67\linewidth]{figures/Fig_05_load_specific_match.pdf}
\caption{Relative magnitude response of T to P at two common shunt loads. The ordinate is $20\log_{10}|H_T/H_P|$; zero indicates equal magnitudes. Lines are analytical responses; open circles are archived LTspice samples. Complete-function equality at 10~k$\Omega$ is established algebraically in the text.}
\label{fig:load-specific}
\end{figure}


\subsection{A bridging capacitor acts on both differential-output gains}

We connected a capacitor \(C_d\) between the two output nodes while retaining the ideal sources and the specified shunt loads. Let \(Z'_P\text{ and }Z'_N\) denote the branch output impedances in parallel with the respective shunt loads, and let \(H_P\) and \(H_N\) denote the loaded transfers before \(C_d\) was added. Solving the two output-node equations gave

\begin{equation}\label{eq:12}
\begin{aligned}
G_{DC}&=\frac{H_P-H_N}{D_d},\\
G_{DD}&=\frac{(H_P+H_N)/2}{D_d},\\
D_d&=1+sC_d(Z'_P+Z'_N).
\end{aligned}
\end{equation}

The shared factor \(D_d\) cancelled from \(G_{DD}/G_{DC}\) wherever the ratio was defined. The bridging capacitor changed both differential-output gains by the same factor and preserved \(G_{DD}/G_{DC}\). Equal loaded branch transfers also preserved \(G_{DC}\) = 0 after the capacitor was connected.

For P/Q with \(R_L=10\,\mathrm{k\Omega}\) and \(C_d=820\,\mathrm{pF}\), the maximum complex-ratio difference was \(5.14\times10^{-12}\) relative to the ratio without the bridging capacitor, evaluated over 145 frequency points. The ratio followed the cancellation of the shared factor in Eq.~\eqref{eq:12}.


Figure~\ref{fig:bridge-gains} displays the changes in the separate gains, while Figure~\ref{fig:bridge-cmrr} compares the corresponding rejection ratios.

\begin{figure}[htbp]
\centering
\includegraphics[width=0.67\linewidth]{figures/Fig_06_bridge_gains.pdf}
\caption{Analytical differential-output gains of the loaded P/Q pair with and without a bridging capacitor. DM denotes unit differential input; CM$\rightarrow$DM denotes unit common-mode input and differential output. Open markers identify analytical samples on the 820~pF curves. Both shunt loads are 10~k$\Omega$.}
\label{fig:bridge-gains}
\end{figure}


\begin{figure}[htbp]
\centering
\includegraphics[width=0.67\linewidth]{figures/Fig_07_bridge_cmrr.pdf}
\caption{Common-mode rejection ratio of the loaded P/Q pair. The analytical curves for $C_d=0$ and 820~pF coincide. Open circles are calculated from archived common-mode and differential-mode LTspice output differences with $C_d=820$~pF. Both shunt loads are 10~k$\Omega$; CMRR is $20\log_{10}|G_{DD}/G_{DC}|$. Absolute gains accompany the ratio in Figure~\ref{fig:bridge-gains}; the independent-branch termination sweep is shown in Figure~\ref{fig:terminal-sweep}.}
\label{fig:bridge-cmrr}
\end{figure}


\FloatBarrier
